\documentclass[10pt,a4paper]{amsart}
\usepackage[margin=19mm]{geometry}
\usepackage{amsmath,amssymb,mathtools}
\usepackage[scr=rsfs,cal=euler]{mathalfa}
\usepackage{xfrac}
\usepackage[unicode,hidelinks,bookmarks=false]{hyperref}
\newcommand{\ie}{i.e.\ }
\newcommand{\cf}{cf.\ }
\newcommand{\resp}{respectively\ }
\newcommand{\Subsection}{\subsection}
\providecommand{\nobreakdash}{\nobreak\hbox{-}}
\setlength{\emergencystretch}{2em}
\multlinegap0pt
\usepackage{amssymb}
\usepackage{enumerate}
\usepackage{caption}
\newtheorem{theorem}{Theorem}
\newcommand\Aut{\mathrm{Aut}}
\newcommand\B{\mathbb{B}}
\newcommand\C{\mathbb{C}}
\title{Approximation of biholomorphic maps between Runge domains by holomorphic automorphisms}
\author{Franc Forstneric}
\date{Source: arXiv:2505.15129v2 (CC BY 4.0)}
\begin{document}
\maketitle

\noindent\emph{Source and licence.} Approximation of biholomorphic maps between Runge domains by holomorphic automorphisms. Version arXiv:2505.15129v2 (CC BY 4.0), distributed by arXiv under the Creative Commons Attribution 4.0 International licence, http://creativecommons.org/licenses/by/4.0/, as recorded in the arXiv licence metadata for this exact version and shown on the abstract page https://arxiv.org/abs/2505.15129v2. Full licence notice: \texttt{licenses/arxiv-2505.15129-CC-BY-4.0.txt}.\par\smallskip

\noindent\textit{Editorial scope.} Counted unit: the article's theorem th:Runge (source0.tex lines 422--430). The article's complete original proof of it is delivered with it (source0.tex lines 467--571, one \texttt{proof} environment of 105 lines, which the article places away from the statement). No mathematics is written, repaired or re-derived: the statement and the proof are the article's own lines, in the article's own notation, and no retained line is substituted or re-flowed.  No further source-local context is needed: every cross-reference of every retained line resolves inside the two delivered spans.  Nothing was omitted from any retained span, so the package carries no omission record.  The retained lines cite 3 works, whose entries are transcribed verbatim from the article's own bibliography source in the e-print and delivered with short editorial labels recorded in the item's reference mapping.  No retained line contains a figure environment, an image-inclusion command or drawing code, so no image is delivered and the package declares no asset dependency.  Editorial formatting only, in addition: an AMS \texttt{amsart} preamble that restates the article's own declarations and macros used by the retained lines, namely the environments \texttt{theorem} on separate counters, together with the standard packages those lines need.  The retained environments print in this document as Theorem 1 in order of appearance, whereas the article prints the same items with the numbering of its own full document.  The complete native source of the article as distributed in the arXiv e-print for this exact version is bundled unchanged as \texttt{sources/178482/source0.tex}.
\begin{theorem}\label{th:Runge}
Assume that $V$ is a complete holomorphic vector field 
on $\C^n$, $n>1$, with a globally attracting 
fixed point $0\in \C^n$ and the domain $0\in \Omega\subset \C^n$
is invariant under the positive time flow of $V$.
Then, every biholomorphic map from $\Omega$ 
onto a Runge domain in $\C^n$ can be approximated uniformly 
on compacts in $\Omega$ by holomorphic automorphisms of $\C^n$.
\end{theorem}

\begin{proof}[Proof of Theorem \ref{th:Runge}] 
Let $V$ and $\Omega\subset \C^n$ be as in the theorem.
By \cite[Theorem 1.1]{ChatterjeeGorai2025X}, 
$\Omega$ is Runge in $\C^n$. We shall prove 
that every biholomorphic map $F:\Omega\to \Omega'$ 
onto a Runge domain $\Omega'\subset\C^n$ 
is a limit of holomorphic automorphisms
of $\C^n$, uniformly on compacts in $\Omega$. 

We may assume that $F(0)=0$ and the derivative
$DF(0)$ is the identity map. 
Thus, $F$ is a small perturbation of the identity near the origin.
In particular, choosing $\epsilon >0$ small enough, 
we have that $\B(0,\epsilon)\subset\Omega$ and 
the image $F(\B(0,\epsilon))$ is convex. Hence, the restricted map
$F:\B(0,\epsilon)\to F(\B(0,\epsilon))$ is a limit of  
holomorphic automorphisms by \cite[Theorem 2.1]{AndersenLempert1992}.
Fix such an $\epsilon$. Note that for each $t\ge 0$ the domain 
$\Omega_t:=\phi_t(\Omega)\subset \Omega$
is Runge in $\C^n$ (and hence in $\Omega$) since $\phi_t\in\Aut(\C^n)$. 

Assume first that $\overline \Omega$ is compact.  
Conditions (1) and (2) on the vector field 
$V$ imply that there is a $t_0>0$ such that 
$\phi_{t}(\overline\Omega)\subset \B(0,\epsilon)$ for all $t\ge t_0$.
Indeed, given a point $p\in \overline \Omega$, condition (1) 
gives a number $t(p)>0$ such that $\phi_{t(p)}(p)\in \B(0,\delta)$.
By continuity, there is a neighbourhood $U_p\subset \C^n$
of $p$ such that $\phi_{t(p)}(U_p)\subset \B(0,\delta)$.
This gives a finite open cover $U_1,\ldots,U_m$ of $\overline \Omega$
and numbers $t_1>0,\ldots,t_m>0$ such that 
\begin{equation}\label{eq:delta}
	\phi_{t_j}(U_j)\subset \B(0,\delta)\ \ \text{holds for $j=1,\ldots,m$}.
\end{equation}
Set $t_0=\max\{t_1,\ldots,t_m\}$. By property (2) of the flow 
and \eqref{eq:delta} we have that 
$\phi_t(\overline \Omega)\subset \B(0,\epsilon)$ for all $t\ge t_0$,
which proves the claim.

Recall that $\Omega'=F(\Omega)$. 
Let $\psi_t:\Omega'\to \Omega'$ for $t\ge 0$ be the unique 
holomorphic map which is $F$-conjugate to $\phi_t:\Omega\to\Omega$,
defined by the condition
\[
	F\circ \phi_t = \psi_t \circ F\ \ \text{for all $t\ge 0$}.
\]
Thus, $\psi_t$ maps $\Omega'$ biholomorphically onto the domain
$
	\Omega'_t:=\psi_t(\Omega') =F(\Omega_t)\subset\Omega'
$
for every $t\ge 0$, and $\psi_0$ is the identity on $\Omega'$. 
Since $\Omega_t$ is Runge in $\Omega$ for every $t\ge 0$ 
and the map $F:\Omega\to\Omega'$ is biholomorphic, we infer
that $\Omega'_t$ is Runge in $\Omega'$, 
and hence also in $\C^n$ (since $\Omega'$ is Runge in $\C^n$). 
Consider the family of maps 
\begin{equation}\label{eq:isotopy}
	F_t = F\circ \phi_t : 
	\Omega \stackrel{\cong}{\longrightarrow} \Omega'_t,
	\quad t\ge 0.
\end{equation}
This is an isotopy of biholomorphic maps from the Runge domain $\Omega$
onto the family of Runge domains $\Omega'_t\subset \C^n$, 
with $F_t$ depending smoothly on $t$. Since 
$\overline \Omega_{t_0}= \phi_{t_0}(\overline \Omega)\subset \B(0,\epsilon)$, 
the restricted map $F:\B(0,\epsilon)\to \C^n$ is a limit of 
automorphisms of $\C^n$ and $\phi_{t_0}\in\Aut(\C^n)$,
the map $F_{t_0}=F\circ\phi_{t_0}$ is 
also a limit of automorphisms of $\C^n$. By 
\cite[Theorem 1.1]{ForstnericRosay1993} 
it follows that every map $F_t$ in the isotopy \eqref{eq:isotopy}
is a limit of automorphisms of $\C^n$. In particular, this holds for
the map $F_0=F:\Omega\to\Omega'$.

This proves the theorem in the case when $\overline\Omega$
is compact. The general case follows by observing that 
$\Omega$ is exhausted by relatively compact 
domains $\Omega_0\Subset \Omega$ containing the origin
which are invariant under the positive time flow of $V$. 
To see this, choose an open relatively compact subset $W$ of $\Omega$ 
and set $\Omega_0=\bigcup_{t\ge 0} \phi_t(W)\subset \Omega$.
Obviously, $\Omega_0$ is open and positive time invariant.
Pick $\epsilon>0$ such that $\overline{\B(0,\epsilon)}\subset\Omega$.
We see as before that there is a number $t_0>0$ 
such that $\phi_t(\overline W)\subset \B(0,\epsilon)$ for all $t\ge t_0$.
It follows that 
\[
	\Omega_0 \subset 
	\bigcup_{0\le t\le t_0}\phi_t(\overline W) \cup 
	\overline{\B(0,\epsilon)}.
\]
Since the first set on the right hand side is compact
and contained in $\Omega$, we see that 
$\overline \Omega_0\subset\Omega$.
By \cite[Theorem 1.1]{ChatterjeeGorai2025X}, 
$\Omega_0$ is Runge in $\C^n$, and hence in $\Omega$.
If follows that $F(\Omega_0)=\Omega'_0$ is Runge 
in $\Omega'=F(\Omega)$, and hence also in $\C^n$ 
(since $\Omega'$ is Runge in $\C^n$). The above argument 
in the special case then shows that $F:\Omega_0\to  \Omega'_0$ 
is a limit of holomorphic automorphisms of $\C^n$ uniformly on
compacts in $\Omega$. 
By the construction, $\Omega_0$ can be chosen to contain 
any given compact subset of $\Omega$, which proves the theorem.
\end{proof}

\begin{thebibliography}{99}
\bibitem[Anders]{AndersenLempert1992} E.~Anders{\'e}n and L.~Lempert. \newblock On the group of holomorphic automorphisms of {${\mathbb C}^n$}. \newblock {\em Invent. Math.}, 110(2):371--388, 1992.
\bibitem[Chatte]{ChatterjeeGorai2025X} S.~Chatterjee and S.~Gorai. \newblock Approximations on certain domains of $\mathbb{C}^{n}$. \newblock arXiv e-prints, 2025, \url{https://arxiv.org/abs/2208.11107}.
\bibitem[Forstn]{ForstnericRosay1993} F.~Forstneri\v{c} and J.-P. Rosay. \newblock Approximation of biholomorphic mappings by automorphisms of {${\mathbb C}^n$}. \newblock {\em Invent. Math.}, 112(2):323--349, 1993. \newblock {Erratum: {\it Invent. Math.}, 118(3):573--574, 1994}.
\end{thebibliography}
\end{document}
